\section{Conclusiones del análisis de la unidad}
Las cinco conclusiones desarrollan aspectos distintos de la práctica. Se dejan como borrador colectivo para lectura y ajustes de los tres integrantes antes de la entrega; no se asigna a ninguna persona una conclusión que no haya revisado. Se conserva la cantidad de cinco solicitada por la guía, aunque el equipo real esté formado por tres miembros.

\subsection{El problema comercial determina el alcance}
La descripción inicial del PFC muestra que un proyecto de ventas no puede partir solamente del deseo de incorporar un componente de inteligencia artificial. El equipo identificó un problema de facturación y consulta de información de productos que afecta al trabajo diario de caja y de inventario, de acuerdo con la base documental disponible \cite{basePFC2026}. Este hallazgo obliga a considerar primero las tareas de búsqueda, cálculo, control de existencias y recuperación de operaciones. Las recomendaciones tienen sentido si acompañan ese proceso y no lo sustituyen. Para nuestra formación, la lección es que una funcionalidad llamativa necesita justificarse por una necesidad del negocio. El siguiente paso será contrastar esa necesidad con las personas que utilizan el establecimiento y registrar sus decisiones.

\subsection{Un requisito necesita condiciones para comprobarse}
Durante la elaboración de la matriz se hizo evidente que escribir «rápido», «seguro» o «fácil de usar» no ofrece una referencia suficiente para evaluar el software. Los requisitos de calidad necesitan una medida, un entorno de prueba, un conjunto de datos y un resultado esperado. Esta idea coincide con los principios de formulación y verificación de ISO/IEC/IEEE 29148:2018 y con el uso del modelo de calidad de ISO/IEC 25010:2023 \cite{iso29148,iso25010}. Los umbrales registrados en el PFC son objetivos académicos y no mediciones obtenidas en producción. Aun así, permiten anticipar cómo se efectuaría una prueba y qué evidencia debería conservarse. La experiencia también muestra que esos valores deben negociarse para no convertir supuestos en obligaciones que el negocio no aprobó.

\subsection{Los actores evitan decisiones técnicas sin contexto}
La revisión de los interesados permitió separar la autoridad de gerencia, la experiencia del vendedor, el control de inventario y las necesidades del comprador. Todos pueden contribuir a los requisitos, pero sus decisiones no tienen el mismo alcance. Por ejemplo, un desarrollador puede explicar cómo se implementa una transacción, aunque no debería definir de forma unilateral la política comercial de descuentos o devoluciones. El análisis de poder e interés sirve para planificar esas conversaciones y decidir a quién consultar sobre cada tema \cite{pmbok2021,pohl2025}. En esta etapa la matriz representa una propuesta del equipo y no una clasificación aprobada por la organización. Su utilidad práctica está en orientar el levantamiento de información y evitar que se dé por válida una regla que solo se dedujo del diagrama de datos.

\subsection{La recomendación con IA exige una decisión humana}
El modelo ER incluye el registro de consultas de recomendación y una caché técnica, lo que ayuda a visualizar qué datos podrían circular por el módulo de IA. Sin embargo, una respuesta almacenada no demuestra que el material sugerido sea eléctricamente adecuado para un uso específico \cite{erPFC2026}. Esto llevó a proponer que el vendedor pueda revisar y descartar la sugerencia, consultar las especificaciones y evitar enviar datos personales innecesarios a un servicio externo. La responsabilidad final sobre el producto ofrecido debe mantenerse en el proceso comercial. Esta distinción resulta importante tanto para el PFC como para nuestra formación profesional: la automatización aporta información, pero su calidad necesita criterios explícitos, supervisión y mecanismos para manejar errores o indisponibilidad del motor.

\subsection{La trazabilidad debe continuar después del informe}
El ejercicio permitió relacionar el problema del negocio, los interesados, el modelo de datos, los requisitos y los casos de validación que todavía deben ejecutarse. Esa relación debe mantenerse cuando cambien precios, reglas de inventario, interfaces o proveedores. La norma ISO/IEC/IEEE 29148:2018 y el marco de Pohl respaldan la idea de conservar el origen de las necesidades y revisar sus consecuencias en el sistema \cite{iso29148,pohl2025}. En términos prácticos, un cambio de requisito debería indicar qué versión modifica, por qué se decidió y qué comprobación se repetirá. El historial de GitHub ayuda a documentar versiones, pero no reemplaza la aprobación de usuarios reales. Como siguiente paso del equipo queda verificar en campo los supuestos que permanecen abiertos y completar las evidencias que exige la práctica.
